\section*{Chapitre \ref{ch:rs:performance-measurement} --- Mesure de la performance}

\begin{solution}{exo:rs:performance-measurement:1}
Environ $\sqrt{252/1\,260} = 0.45$ (le terme $\mathrm{SR}^2/2$ est négligeable par jour)~; l'intervalle à 95~\% va de 0.12 à 1.88.
\end{solution}

\begin{solution}{exo:rs:performance-measurement:2}
$0.144/0.336 = 0.43$. Le \omterm{def:rs:performance-measurement:drawdown}{drawdown maximal} n'est qu'un tirage dans une distribution large~: le maximum sur cinq ans de la même stratégie dépasse 30~\% dans 14~\% des trajectoires et reste
bien en dessous de ce seuil dans la plupart~; un ratio dont le dénominateur contient un cas extrême hérite de cette instabilité.
\end{solution}

\begin{solution}{exo:rs:performance-measurement:3}
Sur $n$ périodes, $\sum r^+ / \sum r^- = (\frac1n\sum r^+)/(\frac1n\sum r^-) = \mathbb{E}_n[(r - 0)^+]/\mathbb{E}_n[(0 - r)^+]$, le \omterm{def:rs:performance-measurement:sortino}{ratio Oméga} empirique en zéro.
\end{solution}

\begin{solution}{exo:rs:performance-measurement:4}
$\sum\theta_j^2 = 0.25 + 0.09 + 0.04 = 0.38$~; facteur de volatilité $\sqrt{0.38} = 0.616$~; ratio de Sharpe annoncé $0.6/0.616 = 0.97$~;
$\rho_1 = (0.5 \times 0.3 + 0.3 \times 0.2)/0.38 = 0.55$, $\rho_2 = 0.5 \times 0.2/0.38 = 0.26$.
\end{solution}

\begin{solution}{exo:rs:performance-measurement:5}
$12/\sqrt{12 + 2(11 \times 0.55 + 10 \times 0.26)} = 2.22$ contre $\sqrt{12} = 3.46$~: l'\omterm{def:rs:performance-measurement:annualisation}{annualisation} naïve surestime le ratio de Sharpe d'un facteur de
$3.46/2.22 = 1.56$, soit 56~\%.
\end{solution}

\begin{solution}{exo:rs:performance-measurement:6}
Le bêta estimé moyenne un profil de gains à rupture~: sur 100 années simulées, le bêta quotidien vaut 0.98 les jours de baisse de l'indice et 0.71 les jours de hausse, et en krach le delta du put
tend vers un, si bien que la perte est bien supérieure à 0.45 fois la baisse de l'indice. L'historique affiché (0.59 et 0.54) ne contient aucun krach pour révéler la rupture.
\end{solution}

\begin{solution}{exo:rs:performance-measurement:7}
Sur un an~: 0.6~\% pour la stratégie de tendance, 11~\% pour la volatilité courte~; sur dix ans~: 31~\% et 65~\%. Les pertes de la volatilité courte arrivent par sauts
(un mois de krach emporte une large part du capital d'un coup), si bien qu'une baisse de 30~\% n'a pas besoin de l'accumulation lente de mauvais jours qu'exige une diffusion.
\end{solution}

\begin{solution}{exo:rs:performance-measurement:8}
Le \omterm{def:rs:performance-measurement:sortino}{ratio de Sortino} du fonds (1.87 dans l'exemple du chapitre) et son bilan proviennent d'un historique sans krach~; une mesure du risque de baisse ne peut pas voir les baisses
qui manquent à l'échantillon. Son modèle donne une probabilité de 41~\% d'observer un drawdown de 30~\% en cinq ans, contre 14~\% pour le fonds de tendance, et sur 100 ans son asymétrie vaut $-1.23$.
Mesurer l'exposition de la stratégie (un put vendu) et la soumettre à un test de résistance, et non à l'historique.
\end{solution}

\begin{solution}{pb:rs:performance-measurement:1}
\begin{enumerate}
\item 1.00 et 1.09, chacun avec une erreur type iid de 0.45 (HAC 0.46 et 0.45).
\item L'erreur type dépend de la longueur de l'historique en années, environ $\sqrt{q/T}$, et les quatre couvrent cinq ans.
\item 0.63 contre 0.46~: des rendements autocorrélés portent moins d'information que leur nombre ne le laisse penser.
\item 4.04 avec une erreur type de 0.45 (HAC 0.52)~: un edge réel, mais rien sur la capacité, sur l'encombrement (crowding) ni sur ce qu'une journée défavorable peut coûter à grande échelle.
\item 7.4~\% annoncé, 11.2~\% réel.
\item 0.99 annoncé, 0.61 avec la correction de Lo, 0.59 une fois le lissage corrigé, 0.66 réel.
\item $(0.43, 0.36, 0.21)$, à partir des premières autocorrélations de 0.64 et 0.25.
\item Ses rendements sont autocorrélés, si bien que la volatilité de la somme annuelle est supérieure à $\sqrt{12}$ fois la volatilité mensuelle.
\item Tendance $-33.6~\%$ et 791 jours~; volatilité courte $-14.8~\%$ et 48 jours~; rendements lissés $-11.0~\%$ et 22 mois~; teneur de marché $-4.2~\%$ et 72 jours.
\item Pour la volatilité courte, l'historique est meilleur sur les trois ratios (Calmar 0.78, Sortino 1.87, Oméga 1.49 contre 0.43, 1.48, 1.17)~; elle a aussi la queue la plus lourde.
\item $+5.38$ et 102.5 dans l'historique~; $-1.23$ et 121 sur 100 ans.
\item 9~\%.
\item 48 jours de négociation~; une rotation unilatérale de 1.86~\% du capital par jour.
\item \textbf{Résultat nommé.} Un drawdown de 30~\% dans les cinq ans~: 14~\% pour la stratégie de tendance, 41~\% pour la volatilité courte. Le ratio de Sharpe du portefeuille lissé a une
erreur type de 0.63 en tenant compte de son autocorrélation (HAC), contre 0.46 sous l'hypothèse d'indépendance.
\item Tendance~: 0.83 et $-37~\%$. Volatilité courte~: 0.62 et $-72~\%$.
\item 0.45 (erreur type 0.09) dans l'historique~; bêtas des jours de baisse et des jours de hausse de 0.59 et 0.54 dans l'historique, 0.98 et 0.71 sur 100 ans.
\item La stratégie de tendance, ou la stratégie de volatilité courte dimensionnée seulement sur son krach (une perte en scénario de stress comprise dans le budget de drawdown), et non sur la volatilité de son historique.
\item Son erreur type (iid et HAC), la longueur de l'historique, le \omterm{def:rs:the-research-process:log}{nombre d'essais} qui ont précédé, et la probabilité de drawdown du modèle de la stratégie.
\item Vendre des options ou d'autres risques de queue, lisser les valorisations, ou choisir la fenêtre de publication~; une mesure insensible à la manipulation ou un test de résistance répond à chacun.
\item Que le ratio de Sharpe réel de la stratégie se situe probablement entre 0.1 et 1.9, et rien sur sa queue.
\end{enumerate}
\end{solution}

\begin{solution}{iq:rs:performance-measurement:1}
L'erreur type vaut environ $\sqrt{252/756} = 0.58$, si bien que l'estimation se situe à 1.7 erreurs types de zéro~: on peut être assez confiant, sans l'être beaucoup, avant toute correction
pour le nombre de stratégies testées (chapitre~20).
\end{solution}

\begin{solution}{iq:rs:performance-measurement:2}
Vendre des options hors de la monnaie ou d'autres risques de queue (gains réguliers, pertes rares et importantes), lisser les valorisations des positions peu liquides, ou choisir la fenêtre. Examiner la
forme sur un long historique simulé ou soumis à un stress, le \omterm{def:rs:performance-measurement:beta}{bêta des marchés} baissiers, l'autocorrélation des rendements et les positions elles-mêmes~; Goetzmann et ses coauteurs
proposent une mesure qui ne peut pas être manipulée.
\end{solution}

\begin{solution}{iq:rs:performance-measurement:3}
Elle gonfle le ratio de Sharpe annualisé (des valorisations obsolètes abaissent la volatilité mesurée) et le rend moins précis. Il faut corriger
l'\omterm{def:rs:performance-measurement:annualisation}{annualisation} de Lo corrigée de l'autocorrélation, retirer le \omterm{def:rs:performance-measurement:smoothing}{lissage des rendements} avec le profil estimé de Getmansky, Lo et Makarov, et utiliser les erreurs types HAC.
\end{solution}

\begin{solution}{iq:rs:performance-measurement:4}
Pour un mouvement brownien à dérive positive, il croît comme le logarithme de l'horizon (ou comme une racine carrée sans dérive). Le maximum d'un \omterm{def:rs:vectorised-backtests:backtest}{backtest} n'est qu'un seul tirage, et
il ne peut pas contenir les événements qui manquent à l'échantillon~; il faut utiliser la distribution des drawdowns sous un modèle et des tests de résistance.
\end{solution}

\begin{solution}{iq:rs:performance-measurement:5}
La distribution des pertes de ses mauvaises journées (les 1.5~\% de journées de sélection adverse), la capacité et la manière dont l'edge varie avec la taille, ses limites d'inventaire, la corrélation de ses
pertes avec le stress de marché, et si le P\&L se décompose en capture du spread et en sélection adverse qui ont un sens (chapitre~23).
\end{solution}

\begin{solution}{iq:rs:performance-measurement:6}
Avec $\hat\mu$ et $\hat\sigma^2$ indépendants pour des données gaussiennes, $\operatorname{Var}\hat\mu = \sigma^2/T$ et $\operatorname{Var}\hat\sigma^2 = 2\sigma^4/T$. La
méthode delta sur $\mathrm{SR} = \mu/\sigma$, dont les dérivées sont $1/\sigma$ et $-\mu/(2\sigma^3)$, donne $\operatorname{Var}\widehat{\mathrm{SR}} = 1/T + \mu^2/(2\sigma^2 T) =
(1 + \mathrm{SR}^2/2)/T$.
\end{solution}
